\begin{frame}{Критерий качества и методы оптимизации}
	\vfill % Автоматический баланс отступов сверху и снизу слайда
	
	Целевой функционал:
	\begin{tcolorbox}[
		colframe=Burgundy!35,       % Брусничный контур
		colback=Burgundy!1!Ivory,    % Мягкий брусничный отлив
		coltext=DeepBlue,            % Синий цвет для текста внутри
		boxrule=0.5pt, arc=3pt, 
		top=3pt, bottom=3pt, left=4pt, right=4pt,
		before=\vspace{2pt}, after=\vspace{3pt}
		]
		\centering
		$\displaystyle
		\theta^* = \arg\min_{\theta \in \Omega} \mathcal{F}(\theta) \quad
		\mathcal{F}(\theta) = \frac{1}{M} \sum_{j=1}^{M} \frac{\|E_j(\theta)\|_\infty}{\|\bar E_j\|_\infty}$
	\end{tcolorbox}
	\vspace{-2pt}
	{\footnotesize $E_j(\theta)$ — вектор погрешности на $j$-й тестовой задаче, \\ $\bar E_j$ — вектор погрешности классического метода (непараметрического 6-го порядка).}
	
	\vspace{0.05cm}
	
	\begin{columns}[T]
		
		% --- ЛЕВАЯ КОЛОНКА: СЕТОЧНЫЙ ПОИСК ---
		\begin{column}{0.45\textwidth}
			Сеточный поиск:
			\begin{tcolorbox}[
				colframe=DeepBlue!35, 
				colback=DeepBlue!1!Ivory, 
				coltext=DeepBlue, 
				boxrule=0.5pt, arc=3pt, 
				top=3pt, bottom=3pt, left=4pt, right=4pt,
				before=\vspace{2pt}, after=\vspace{0pt},
				equal height group=opt_methods, % <--- 1. Объединяем рамки в одну группу по высоте
				valign=center                  % <--- 2. Центрируем разреженный текст по вертикали
				]
				\centering
				$\displaystyle \theta \in [0, 1]^4 \cap \Omega$
				
				\vspace{0.25cm}
				\footnotesize{$\displaystyle \Delta c_i: \: 0.1 \rightarrow 0.04 \rightarrow 0.02 \rightarrow 0.01$}
			\end{tcolorbox}
		\end{column}
		
		% --- ПРАВАЯ КОЛОНКА: МЕТОД РОЯ ЧАСТИЦ (PSO) ---
		\begin{column}{0.55\textwidth}
			Метод роя частиц (PSO):
			\begin{tcolorbox}[
				colframe=DeepBlue!35, 
				colback=DeepBlue!1!Ivory, 
				coltext=DeepBlue, 
				boxrule=0.5pt, arc=3pt, 
				top=3pt, bottom=3pt, left=4pt, right=4pt,
				before=\vspace{2pt}, after=\vspace{0pt},
				equal height group=opt_methods  % <--- 1. Тот же идентификатор группы
				]
				{\footnotesize Коррекция координат частиц:}
				{\centering $\displaystyle x_i^{(k+1)} = x_i^{(k)} + v_i^{(k+1)}$ \par}
				
				\vspace{0.15cm}
				{\footnotesize Модификация вектора скорости:}
				{\centering $\displaystyle
					\begin{aligned}
						v_i^{(k+1)} = \omega v_i^{(k)} &+ c_1 r_1 \left( p_{i,\mathrm{best}} - x_i^{(k)} \right) + \\
						&+ c_2 r_2 \left( p_{\mathrm{neigh}} - x_i^{(k)} \right)
					\end{aligned}$ \par}
			\end{tcolorbox}
		\end{column}
		
	\end{columns}
	
	\vfill
\end{frame}